\begin{frame}{Ratios de rentabilidad}
\color{rentabilidad}\textbf{Capacidad de generar beneficios respecto a ventas, activos o patrimonio}\color{black}
\end{frame}

\begin{frame}{Margen neto}
\[
\text{Margen neto}=\frac{\text{Utilidad neta}}{\text{Ventas netas}}\times 100
\]
Mide cuánto queda como utilidad neta por cada unidad monetaria vendida.
\end{frame}

\begin{frame}{Caso práctico: margen neto}
\casoeducativo
\[
\frac{135,000.00}{1,500,000.00}\times 100=9.00\text{ \char37}
\]
\textbf{Interpretación:} por cada \soles{100.00} de ventas, la empresa obtiene \soles{9.00} de utilidad neta.

\textbf{Decisión:} revisar precios, costos y gastos operativos.
\end{frame}

\begin{frame}{Rentabilidad sobre activos: ROA}
\[
\text{ROA}=\frac{\text{Utilidad neta}}{\text{Activo promedio}}\times 100
\]
\[
\frac{135,000.00}{1,500,000.00}\times 100=9.00\text{ \char37}
\]
\textbf{Interpretación:} los activos generaron una rentabilidad neta de 9.00 \char37{}.
\end{frame}

\begin{frame}{Rentabilidad sobre patrimonio: ROE}
\[
\text{ROE}=\frac{\text{Utilidad neta}}{\text{Patrimonio promedio}}\times 100
\]
\[
\frac{135,000.00}{900,000.00}\times 100=15.00\text{ \char37}
\]
\textbf{Interpretación:} el patrimonio generó una rentabilidad de 15.00 \char37{}.
\end{frame}

\begin{frame}{Comparación ROA y ROE}
\begin{itemize}
  \item ROA evalúa eficiencia de los activos.
  \item ROE evalúa rentabilidad para propietarios.
  \item Una diferencia amplia puede relacionarse con apalancamiento.
  \item La lectura requiere revisar deuda y costo financiero.
\end{itemize}
\end{frame}
